\chapter{Introduction}
\section{Motivation and problem statement}
A useful document repository must do more than retain bytes. Users need to organise material, grant limited access and recover information when the exact filename or wording is unknown. A conventional file browser provides an explicit hierarchy, whereas document questions often concern relationships across passages and files. Combining these behaviours raises a practical engineering problem: the same access rules must govern a download, a search result, a retrieved context passage and a saved conversation.

NimbusVault addresses this problem through a storage application with a private document-intelligence subsystem. PostgreSQL remains the authority for application identity and permissions; S3 holds original objects. The intelligence service receives an authorised file scope from the backend rather than deciding that scope from a natural-language request. A retrieved passage is therefore not automatically permission to expose its parent document.

The original October proposal described a scalable cloud-storage architecture and an intended RAG layer. The implemented system retains that separation of concerns while making several choices more concrete: Express rather than NestJS, SQLite-backed private retrieval rather than a mandatory managed vector database, and self-hosted Qwen inference rather than a paid language-model API. Complete EC2 hosting is prepared but not yet accepted. These distinctions prevent the proposal's intentions from being mistaken for completed deployment evidence.
\section{Objectives and realised scope}
Table~\ref{tab:objectives} relates the initial objectives to observed outcomes. The evaluated unit is a bounded academic prototype with real integrations, not an independently audited commercial storage service.
\par\begin{table}[htbp]\caption{Original objectives and evidence-backed outcomes}\label{tab:objectives}\small\begin{tabularx}{\linewidth}{p{.30\linewidth}YY}\toprule Objective&Implemented outcome&Qualification\\\midrule
Secure cloud originals&Private S3 and authenticated adapters&Historical synthetic upload and hash readback\\
Structured organisation&Folders, ownership, move, trash and restore&PostgreSQL-backed application tests\\
Controlled sharing&User permissions and bounded share links&Revocation and inheritance tested\\
Semantic discovery&MiniLM, BM25 and hybrid retrieval&84-query development regression\\
Natural-language answers&Qwen-generated multi-document answers&40 strict local Linux cases\\
Cloud deployment&Native Linux packaging and safeguards&EC2 provisioning/acceptance pending\\
High availability&Separation of service responsibilities&No distributed-availability demonstration\\\bottomrule\end{tabularx}\end{table}\par
\section{Contributions and report organisation}
The principal engineering contributions are permission-preserving integration, evidence-grounded answer handling, a measured hybrid retriever, donor-independent model packaging and a cautious distinction between cloud processing and private AI computation. Later chapters examine each contribution through source structure and identifiable evidence. Requirements and architecture precede implementation; the evaluation then separates observed results from limitations, costs and future deployment work.
\screen{public-landing}{The implemented public frontend, rendered from the verified static build. This interface view is not evidence of a running EC2 backend.}{landing}
\FloatBarrier
